\chapter{Evidence}\label{r03:app:evidence}

Every number in this report comes from one of three sources:
\begin{enumerate}
  \item the raw logs written on the board, copied unchanged to \file{docs/notes3/board-logs/}
        (newer than \rcommitws; the folder's \file{README.md} lists them);
  \item runs made while writing this report, on the committed code, marked ``repeated for this
        report'' in the text;
  \item the working log of the project, \ghws{docs/notes3/rvv-8x-speedup-progress-fable5.2.md},
        for the instruction counts read from cross-compiled assembly on the laptop. It records
        the work day by day, in numbered work steps. This report does not use those step numbers
        and can be read without the log.
\end{enumerate}

\paragraph{Board clock.} The board's clock runs about 1.7~days behind the laptop's. All times in
the board logs, and the start times in \cref{r03:tab:runs}, are board-clock times (2026-11-12);
the same work happened on 2026-11-13 and 2026-11-14 by the laptop's clock. The logs were copied
at board-clock 2026-11-12 20:08, except \file{e2e-cg-gcc-f32.log} (20:29),
\file{e2e-cg-gcc-f64.log} and \file{e2e-halo-stage.log} (21:26), which were copied after they
were complete.

\section*{Session scripts}

\begin{table}[htbp]
  \centering\small
  \caption{The session scripts in \texttt{misc/x60-lab/}, what they measured, and where the results
    are used. Runs that are not listed were started by hand with
    \texttt{scripts/walberla\_lbm\_bench.sh} or \texttt{misc/x60-lab/lab/run.sh}.}\label{r03:tab:sessions}
  \begin{tabularx}{\linewidth}{@{}lXl@{}}
    \toprule
    \textbf{Script} & \textbf{Measures} & \textbf{Used in}\\
    \midrule
    \file{queue_aa_tests.sh} & hand AA kernels (f32); \ps{} RVV tests with prefetch &
      \ref{r03:sec:lever-aa}, \ref{r03:sec:ver-ps}\\
    \file{queue_aa64.sh} & hand AA kernels (f64), hardware counters & \ref{r03:sec:lever-aa}\\
    \file{queue_lab2.sh} & generated AA kernels, LMUL, Clang's X60 model &
      \ref{r03:sec:lever-aa}, \ref{r03:sec:lever-sched}\\
    \file{queue_lab3.sh}, \file{queue_vf2_gate.sh} & new instruction selection: acceptance test,
      \ps{} tests, comparison & \ref{r03:sec:lever-selector}, \ref{r03:sec:ver-ps}\\
    \file{queue_e2e_aa.sh} & run B; 4 and 8 processes & \ref{r03:sec:res-aa},
      \ref{r03:sec:res-mpi}\\
    \file{queue_lab4.sh} & compilers and scheduling models; lab against application &
      \ref{r03:sec:lever-sched}, \ref{r03:sec:res-labapp}\\
    \file{queue_perf_e2e.sh} & profiles of the AA and pull applications & \ref{r03:sec:res-labapp}\\
    \file{queue_lab5.sh} & row pointers and constants as source edits; alignment, huge pages &
      \ref{r03:sec:lever-rp}, \ref{r03:sec:lever-layout}\\
    \file{queue_e2e_cg.sh} & \ps{} tests and run C & \ref{r03:sec:ver-ps}, \ref{r03:sec:res-cg}\\
    \file{queue_e2e_halo.sh} & keep rule for run C's change; slot-by-slot halo exchange end to end
      & \ref{r03:sec:res-cg}, pending (\ref{r03:sec:open-jobs})\\
    \file{queue_lab6.sh}, \file{queue_lab7.sh}, \file{queue_e2e_final.sh} & pending measurements
      & \ref{r03:sec:open-jobs}\\
    \bottomrule
  \end{tabularx}
\end{table}

\section*{Board logs}
\begin{xltabular}{\linewidth}{@{}>{\raggedright\arraybackslash}p{6.4cm}>{\raggedright\arraybackslash}X@{}}
  \toprule
  \textbf{File in \texttt{docs/notes3/board-logs/}} & \textbf{Used for}\\
  \midrule
  \endhead
  \file{ann-explicit_cse-f32.txt}, \file{pmu-gcc-f32-explicit_cse.tsv},
    \file{pmu-gcc-f32-scalar_cse.tsv} & profile and hardware counters of the original program,
    \cref{r03:sec:cycles}\\
  \file{p2-stream-sweep.txt}, \file{p2-stream-levers.txt}, \file{p2-prefetch.txt} & memory
    probe, \cref{r03:sec:mem}\\
  \file{p2-ops.txt} & instruction probe, \cref{r03:sec:ops}\\
  \file{lab/results.tsv} & every lab run made with \file{run.sh}, one line each
    (\cref{r03:tab:pf} and the hand split kernel, \cref{r03:sec:lever-split})\\
  \file{pmu-lab-*.tsv} & hardware counters of lab kernels, \cref{r03:sec:lever-split,r03:sec:lever-aa}\\
  \file{e2e-gcc-f32.log}, \file{e2e-gcc-f64.log}, \file{e2e-clang-f32.log},
    \file{e2e-clang-f64.log} & run A, \cref{r03:tab:res-pf}\\
  \file{pytest-rvv.log} & \ps{} tests on the board with prefetch, \cref{r03:tab:ver-ps}\\
  \file{lab-aa.log}, \file{lab-aa64.log}, \file{pmu-aa.log} & hand AA kernels,
    \cref{r03:sec:lever-aa}\\
  \file{lab2.log} & \cref{r03:tab:genaa}; LMUL and scheduling, \cref{r03:sec:lever-sched}\\
  \file{lab3.log}, \file{lab3-gate.log}, \file{lab3.ok}, \file{pytest-vf2.log} &
    \cref{r03:tab:selector}; acceptance tests, \cref{r03:sec:ver-ps,r03:sec:ver-board}\\
  \file{e2e-aa-stage.log}, \file{e2e-aa-gcc-f32.log}, \file{e2e-aa-gcc-f64.log} &
    run B, \cref{r03:tab:res-aa,r03:tab:waterfall32,r03:tab:waterfall64}\\
  \file{lab4.log} & \cref{r03:tab:tune,r03:tab:labapp}; the hand kernel's 163.9 and 147.8 cycles
    per cell (\cref{r03:fig:ladder})\\
  \file{lab5.log} & \cref{r03:sec:lever-rp}, \cref{r03:tab:align}\\
  \file{perf-e2e-aa.txt}, \file{perf-e2e-pull.txt} & \cref{r03:tab:perf}\\
  \file{e2e-mpi.log}, \file{walberla-mpi-np*-gcc-*.log} & \cref{r03:tab:mpi}\\
  \file{pytest-cg.log} & \ps{} tests on the board with row pointers, 148 passed\\
  \file{e2e-cg-gcc-f32.log}, \file{e2e-cg-gcc-f64.log} & run C, \cref{r03:tab:res-cg}\\
  \file{e2e-halo-stage.log} & the keep rule's decision for run C's change, \cref{r03:sec:res-cg}\\
  \bottomrule
\end{xltabular}

\section*{Runs made for this report}
On the laptop, at \rcommitps{} and \rcommitwl:
\begin{itemize}
  \item \ps{} \file{tests/nbackend} and \file{tests/codegen}: 435 passed, 15 skipped
        (\cref{r03:sec:ver-ps});
  \item waLBerla \code{ctest -R "LbmScenarios|HaloExchange|CommunicationHiding"}: 24 of 24
        passed, with test programs built after the last change to the pack-info header
        (\cref{r03:sec:ver-wl});
  \item the application's checksum matrix of \cref{r03:tab:ver-app}: 60 runs, after regenerating
        the five variants' sources;
  \item cross-checks of the committed sources: the counts in \cref{r03:lst:aa-odd} (9 constants,
        23 row pointers, 23 prefetches, 19 loads, 23 stores in the vector loop); the submodule
        commits recorded in \rcommitws; the branch heads on GitHub (\code{git ls-remote}), which
        match the three commits;
  \item the lab kernel files behind the naming error of \cref{r03:sec:repro-gen}: the number of
        CSE temporaries in each file, and the checksums in \file{lab/results.tsv}.
\end{itemize}
On the board: the compiler versions, core count, governor, frequency and ISA string of
\cref{r03:tab:hw}.
